\section{Forcing über einem abzählbaren Modell}

\subsection{Bedingungen und generische Filter}

Wir entwickeln Mengenforcing so, dass auch ein von außen nicht
fundiertes Modell als Ausgangsmodell zulässig ist. Die Definitionen
und Ranginduktionen werden zunächst \emph{innerhalb von ZFC}
durchgeführt. Danach wird über einem abzählbaren Modell ein äußerer
Quotient gebildet. Eine äußere rekursive Auswertung möglicherweise
nicht fundierter Namen wäre dagegen unzulässig.

\FormulaDefDeltaK[Forcingordnung und Dichtheit]
{\begin{aligned}
 \mathbb P&=(P,\le,1),\qquad p\le1\quad(p\in P),\\
 p\perp q&\ \Longleftrightarrow\ \neg\exists r\ (r\le p\land r\le q),\\
 D\text{ dicht unter }p
 &\ \Longleftrightarrow\ \forall q\le p\ \exists r\le q\ (r\in D)
\end{aligned}}
{B48ForcingOrder}{
  \DeltaRow{Menge}{P\ne\varnothing}
  \DeltaRow{Ordnung}{\le\text{ partiell}}
}
\noindent
\(p\le q\) bedeutet: \(p\) ist die stärkere, informativere Bedingung.
Eine Antikette ist eine Menge paarweise unverträglicher Bedingungen.
Sie heißt maximal in einer Teilmenge \(D\), wenn keine weitere
Bedingung aus \(D\) hinzugefügt werden kann.

\FormulaDefDeltaK[Generischer Filter]
{\begin{aligned}
 &\varnothing\ne G\subseteq P,\\
 &p\in G,\ p\le q\ \Longrightarrow\ q\in G,\\
 &p,q\in G\ \Longrightarrow\ \exists r\in G\ (r\le p,q),\\
 &D\in M,\quad M\models\text{\(D\) dicht in \(\mathbb P\)}
       \ \Longrightarrow\ G\cap D^M\ne\varnothing
\end{aligned}}
{B48GenericDef}{
  \DeltaRow{Modell}{M\models\BXLVIIIZFC,\quad\mathbb P\in M}
}
\noindent
Bei nichttransitivem \(M\) bedeutet \(D^M\) die äußere Menge der
Modellelemente, welche nach der Mitgliedschaftsrelation von \(M\)
in \(D\) liegen. Dasselbe gilt für \(P^M\) und für die Ordnung.
Zur Lesbarkeit werden diese hochgestellten \(M\) im weiteren Text
weggelassen, wenn klar intern gearbeitet wird.

\FormulaThmDeltaK[Existenz generischer Filter]
{\begin{gathered}
 M\models\BXLVIIIZFC,\quad |M|\le\aleph_0,\quad
 M\models p\in\mathbb P\\
 \vdash\ \exists G\quad
 G\text{ ist \(M\)-generisch und }p\in G
 \end{gathered}}
{B48GenericExists}{
  \DeltaRow{Äußere Voraussetzung}{M\text{ abzählbar}}
}
\begin{tabproof}
 \proofstep{1}{M\text{ abzählbar},\quad M\models p\in\mathbb P}{\rA}
 \proofstep{1}{\BXLVIIIText{Zähle von außen alle Elemente
 \(D_0,D_1,\ldots\) von \(M\) auf, welche \(M\) für dichte
 Teilmengen von \(\mathbb P\) hält.}}
 {\BXLVIIIWhy{Teilmenge einer abzählbaren Grundmenge}}
 \proofstep{1}{p_0=p,\qquad p_{n+1}\le p_n,\quad p_{n+1}\in D_n}
 {\BXLVIIIWhy{Interne Dichtheit liefert jeweils ein Modellelement}}
 \proofstep{1}{G=\{q:\exists n<\omega\ (p_n\le q)\}}
 {\BXLVIIIWhy{Äußere Mengenbildung}}
 \proofstep{1}{\BXLVIIIText{Aufwärtsabschluss ist unmittelbar.
 Zwei Elemente werden durch ein \(p_n\) mit genügend großem
 Index gemeinsam verstärkt. Außerdem liegt
 \(p_{n+1}\in G\cap D_n\).}}
 {\BXLVIIIWhy{Absteigende Folge}}
\end{tabproof}
\noindent
Ist \(D\in M\) nur unter \(p\in G\) dicht, so trifft \(G\) ebenfalls
\(D\): Die Menge
\(D\cup\{q:q\perp p\}\) ist dicht in ganz \(\mathbb P\), und ein
Filter, der \(p\) enthält, kann keine zu \(p\) unverträgliche
Bedingung enthalten. Dieselbe Beobachtung gilt für maximale
Antiketten über die Menge ihrer Verstärkungen.

\subsection{Namen und die Forcingrelation}

\FormulaDefDeltaK[Namen und kanonische Namen]
{\begin{aligned}
 N_0&=\varnothing,\qquad
 N_{\alpha+1}=\mathcal P(N_\alpha\times P),\qquad
 N_\lambda=\bigcup_{\alpha<\lambda}N_\alpha,\\
 \tau\text{ ist ein Name}&\ \Longleftrightarrow\
             \exists\alpha\ (\tau\in N_\alpha),\\
 \check x&=\{(\check y,1):y\in x\},\\
 \rho(\tau)&=\sup\{\rho(\sigma)+1:
                         \exists p\ ((\sigma,p)\in\tau)\}
\end{aligned}}
{B48Names}{
  \DeltaRow{Forcingordnung}{\mathbb P}
}
\noindent
Die Namensstufen sind aufsteigend; jeder Name ist eine Menge von
Paaren aus früheren Namen und Bedingungen. Die Rekursionen sind
durch Rang beziehungsweise Fundiertheit gerechtfertigt.
Alle hier auftretenden Namen gehören zunächst zum Ausgangsuniversum.

\FormulaDefDeltaK[Atomare Forcingrelation]
{\begin{aligned}
 p\Vdash\tau\in\sigma\ \Longleftrightarrow\
 &\ \forall q\le p\ \exists r\le q\
 \exists(\eta,s)\in\sigma\
       (r\le s\land r\Vdash\tau=\eta),\\
 p\Vdash\tau\subseteq\sigma\ \Longleftrightarrow\
 &\ \forall(\eta,s)\in\tau\ \forall q\le p,s\
                       \exists r\le q\ (r\Vdash\eta\in\sigma),\\
 p\Vdash\tau=\sigma\ \Longleftrightarrow\
 &\ (p\Vdash\tau\subseteq\sigma)
       \land(p\Vdash\sigma\subseteq\tau)
\end{aligned}}
{B48AtomicForcing}{
  \DeltaRow{Namen}{\tau,\sigma,\eta}
  \DeltaRow{Bedingungen}{p,q,r,s}
}
\noindent
Die Teilmengenrelation in dieser Definition ist zunächst eine
Hilfsrelation. Die Rekursion ist wohlfundiert: Ordne die beiden
Namensränge absteigend und vergleiche die geordneten Paare
lexikographisch. Jeder echte rekursive Aufruf ersetzt einen
Namen durch einen Namen kleineren Ranges. Gleichheit wird danach
als Konjunktion der beiden schon bestimmten Inklusionen definiert.

\FormulaDefDeltaK[Forcing für zusammengesetzte Formeln]
{\begin{aligned}
 p\Vdash\varphi\land\psi
 &\ \Longleftrightarrow\
       (p\Vdash\varphi)\land(p\Vdash\psi),\\
 p\Vdash\neg\varphi
 &\ \Longleftrightarrow\
       \neg\exists q\le p\ (q\Vdash\varphi),\\
 p\Vdash\exists x\,\varphi(x)
 &\ \Longleftrightarrow\
       \forall q\le p\ \exists r\le q\
       \exists\tau\text{ Name}\ (r\Vdash\varphi(\tau))
\end{aligned}}
{B48FormulaForcing}{
  \DeltaRow{Formeln}{\varphi,\psi\text{ jeweils fest}}
}
\noindent
Dies ist eine metatheoretische Rekursion über gewöhnliche endliche
Formeln: Für jede feste Formel erhält man eine Mengenlehreformel,
die ihre Forcingrelation definiert. Der Namensquantor ist ein
gewöhnlicher Mengenquantor mit der definierbaren Einschränkung
\enquote{ist ein Name}. Es wird keine einzige interne
Wahrheitsrelation für alle Formeln des Mengenuniversums behauptet.

\FormulaThmDeltaK[Monotonie, Dichtheitslemma und Entscheidung]
{\begin{aligned}
 q\le p,\ p\Vdash\varphi&\ \Longrightarrow\ q\Vdash\varphi,\\
 \{q:q\Vdash\varphi\}\text{ dicht unter }p
       &\ \Longrightarrow\ p\Vdash\varphi,\\
 \{p:p\Vdash\varphi\text{ oder }p\Vdash\neg\varphi\}
       &\text{ ist dicht in }\mathbb P
\end{aligned}}
{B48ForcingDensity}{
  \DeltaRow{Formel}{\varphi\text{ mit Namensparametern}}
}
\begin{tabproof}
 \proofstep{}{\BXLVIIIText{Induktion zuerst über die atomare
 Rangrekursion, dann über die Formel. Eine Verstärkung
 verkleinert jeweils nur den Bereich der Bedingungen, die noch
 betrachtet werden müssen. Dies beweist Monotonie.}}
 {\BXLVIIIWhy{Definitionsklauseln}}
 \proofstep{}{\BXLVIIIText{Für atomare Mitgliedschaft und
 Inklusion lässt sich unter jeder vorgegebenen Bedingung
 zunächst eine forcierende Bedingung wählen und danach der
 dort verlangte Mitgliedschafts- oder Gleichheitszeuge.
 So werden zwei aufeinanderfolgende Dichtheitsforderungen
 zu einer. Gleichheit folgt komponentenweise.}}
 {\BXLVIIIWhy{Atomare Klauseln}}
 \proofstep{}{\BXLVIIIText{Für Konjunktion gilt das Dichtheitslemma
 für beide Teile; für Existenz werden wieder zwei dichte
 Verstärkungsschritte hintereinander ausgeführt.
 Ist hingegen das Forcen von \(\neg\psi\) dicht unter \(p\),
 so kann kein \(q\le p\) die Formel \(\psi\) forcen:
 Eine weitere Verstärkung würde beides forcen.}}
 {\BXLVIIIWhy{Formelinduktion; Negationsklausel}}
 \proofstep{}{\BXLVIIIText{Für ein beliebiges \(p\) gibt es
 entweder ein \(q\le p\) mit \(q\Vdash\varphi\), oder die
 Negationsklausel besagt bereits \(p\Vdash\neg\varphi\).
 Daher ist die Entscheidungsmenge dicht.}}
 {\BXLVIIIWhy{Klassische Fallunterscheidung}}
\end{tabproof}

\FormulaThmDeltaK[Gleichheitsgesetze des Forcings]
{\begin{aligned}
 &1\Vdash\tau=\tau,\\
 &p\Vdash\tau=\sigma\ \Longrightarrow\ p\Vdash\sigma=\tau,\\
 &p\Vdash\tau=\sigma,\ p\Vdash\sigma=\eta
                 \ \Longrightarrow\ p\Vdash\tau=\eta,\\
 &p\Vdash\tau=\sigma,\ p\Vdash\varphi(\tau)
                 \ \Longrightarrow\ p\Vdash\varphi(\sigma)
\end{aligned}}
{B48ForcingEquality}{
  \DeltaRow{Namen}{\tau,\sigma,\eta}
  \DeltaRow{Formel}{\varphi}
}
\begin{tabproof}
 \proofstep{}{\BXLVIIIText{Reflexivität folgt durch Ranginduktion.
 Zu \((\nu,s)\in\tau\) und \(q\le s\) liefert
 \(1\Vdash\nu=\nu\) die Mitgliedschaft
 \(q\Vdash\nu\in\tau\). Beide Inklusionsklauseln für
 \(\tau=\tau\) sind damit erfüllt. Symmetrie ist definitorisch.}}
 {\BXLVIIIWhy{Atomare Rekursion}}
 \proofstep{}{\BXLVIIIText{Für Transitivität und atomare
 Substitution führe eine gemeinsame Induktion über die
 absteigend sortierten Tripel der beteiligten Namensränge.
 Beim Übergang von \(\tau\) nach \(\sigma\) und weiter nach
 \(\eta\) liefert jede verlangte Verstärkung sukzessive
 Mitgliedschaftszeugen aus den beiden folgenden Namen.}}
 {\BXLVIIIWhy{Wohlfundierte lexikographische Rangordnung}}
 \proofstep{}{\BXLVIIIText{Genauer: Zu einem Paar \((u,s)\in\tau\)
 und einer Bedingung unter \(p,s\) wähle aus
 \(\tau\subseteq\sigma\) ein Paar \((v,t)\in\sigma\) mit
 stärkerer Bedingung, die \(u=v\) forciert.
 Aus \(\sigma\subseteq\eta\) wähle danach
 \((w,a)\in\eta\) und eine weitere Verstärkung, die \(v=w\)
 forciert. Die Ränge von \(u,v,w\) sind jeweils kleiner als
 die von \(\tau,\sigma,\eta\); Induktion ergibt \(u=w\).
 Dies beweist \(\tau\subseteq\eta\); die Rückrichtung ist gleich.}}
 {\BXLVIIIWhy{Transitivitätsschritt}}
 \proofstep{}{\BXLVIIIText{Für die Substitution in \(u\in v\)
 ersetze zunächst den Mitgliedschaftszeugen durch den
 gleichgesetzten Namen und verwende die gerade bewiesene
 Transitivität. Wird der rechte Name ersetzt, transportiere
 seinen Zeugen über die erzwungene Inklusion in den neuen
 rechten Namen. Dabei sinkt mindestens einer der drei Ränge;
 die übrigen steigen nicht. Das Dichtheitslemma entfernt die
 jeweils eingeführten Verstärkungen.}}
 {\BXLVIIIWhy{Gemeinsamer atomarer Induktionsschritt}}
 \proofstep{}{\BXLVIIIText{Formelinduktion liefert die allgemeine
 Substitution. Konjunktion und Existenz transportieren die
 Teilformeln und Zeugen. Bei Negation würde eine forcierte
 Gegeninstanz nach Symmetrie und Induktion eine Gegeninstanz
 vor der Substitution liefern, entgegen der Negationsklausel.}}
 {\BXLVIIIWhy{Monotonie und Formeldefinition}}
\end{tabproof}

\subsection{Der generische Quotient und das Wahrheitslemma}

\FormulaDefDeltaK[Generische Erweiterung als Quotient]
{\begin{aligned}
 \tau\sim_G\sigma
 &\ \Longleftrightarrow\
 \exists p\in G\ (M\models p\Vdash\tau=\sigma),\\
 [\tau]_G\ E_G\ [\sigma]_G
 &\ \Longleftrightarrow\
 \exists p\in G\ (M\models p\Vdash\tau\in\sigma),\\
 M[G]&=\bigl(\{\text{\(M\)-Namen}\}/\sim_G,\ E_G\bigr)
\end{aligned}}
{B48GenericQuotient}{
  \DeltaRow{Modell}{M\models\BXLVIIIZFC}
  \DeltaRow{Filter}{G\text{ ist \(M\)-generisch}}
}
\noindent
Die \(M\)-Namen bilden von außen eine Teilmenge der Grundmenge von
\(M\), also eine Menge. Die Gleichheitsgesetze und die gemeinsame
Verstärkung endlich vieler Bedingungen aus \(G\) machen
\(\sim_G\) zu einer Äquivalenzrelation und \(E_G\) unabhängig von
der Repräsentantenwahl. Im Folgenden schreiben wir kurz
\(\tau_G=[\tau]_G\) und innerhalb der Erweiterung wieder \(\in\)
anstelle von \(E_G\).

\FormulaThmDeltaK[Wahrheitslemma]
{M[G]\models\varphi(\tau_{1,G},\ldots,\tau_{n,G})
 \ \Longleftrightarrow\
 \exists p\in G\ (M\models p\Vdash\varphi(\tau_1,\ldots,\tau_n))}
{B48TruthLemma}{
  \DeltaRow{Formel}{\varphi\text{ eine feste gewöhnliche Formel}}
}
\begin{tabproof}
 \proofstep{}{\BXLVIIIText{Induktion über den wirklichen endlichen
 Aufbau von \(\varphi\). Atomare Gleichheit und Mitgliedschaft
 gelten nach Definition des Quotienten.}}
 {\BXLVIIIWhy{Quotientendefinition}}
 \proofstep{}{\BXLVIIIText{Für Konjunktion liefert die
 Induktionsvoraussetzung Bedingungen für beide Teile.
 Eine gemeinsame Verstärkung aus \(G\) forciert beide.
 Umgekehrt forciert eine Konjunktion jeden ihrer Teile.}}
 {\BXLVIIIWhy{Filtereigenschaft und Monotonie}}
 \proofstep{}{\BXLVIIIText{Forciert ein \(p\in G\) die Negation,
 kann keine Bedingung aus \(G\) das Positive forcen:
 Eine gemeinsame Verstärkung widerspräche der Negationsklausel.
 Ist umgekehrt das Positive in \(M[G]\) falsch, so trifft \(G\)
 die in \(M\) definierte dichte Entscheidungsmenge.
 Der getroffene Fall kann nur die Negation sein.}}
 {\FormulaRefAuto{B48ForcingDensity}[thm]}
 \proofstep{}{\BXLVIIIText{Forciert \(p\in G\) einen Existenzsatz,
 trifft \(G\) die unter \(p\) dichte Menge von Bedingungen mit
 einem Namenszeugen. Dieser Zeuge erfüllt nach Induktion die
 Formel in \(M[G]\). Umgekehrt besitzt jeder Zeuge in
 \(M[G]\) einen Namensrepräsentanten; eine dessen Formel
 forcierende Bedingung forciert auch den Existenzsatz.}}
 {\BXLVIIIWhy{Generizität, Existenzklausel und Induktionsvoraussetzung}}
\end{tabproof}

\FormulaThmDeltaK[Mitglieder eines ausgewerteten Namens]
{x\in\sigma_G\ \Longleftrightarrow\
 \exists(\eta,s)\in^M\sigma\ (s\in G\land x=\eta_G)}
{B48NameMembers}{
  \DeltaRow{Element}{x\in M[G]}
  \DeltaRow{Name}{\sigma\in M}
}
\begin{tabproof}
 \proofstep{1}{x=\tau_G\in\sigma_G}{\rA}
 \proofstep{1}{\BXLVIIIText{Wähle \(p\in G\) mit
 \(p\Vdash\tau\in\sigma\). Die atomare Klausel gibt eine
 unter \(p\) dichte Menge von Bedingungen \(r\), die für ein
 \((\eta,s)\in\sigma\) sowohl \(r\le s\) als auch
 \(r\Vdash\tau=\eta\) erfüllen.}}
 {\BXLVIIIWhy{Wahrheitslemma und atomare Definition}}
 \proofstep{1}{\BXLVIIIText{Ein solcher Zeuge \(r\) liegt in \(G\);
 dann liegt auch \(s\) in \(G\) und \(x=\eta_G\).}}
 {\BXLVIIIWhy{Generizität und Aufwärtsabschluss}}
 \proofstep{4}{(\eta,s)\in^M\sigma,\quad s\in G,\quad x=\eta_G}{\rA}
 \proofstep{4}{s\Vdash\eta\in\sigma,\qquad x\in\sigma_G}
 {\BXLVIIIWhy{Reflexivität von Gleichheit; atomare Klausel und
 Wahrheitslemma}}
\end{tabproof}

\noindent
Aus den Definitionsklauseln folgen auch die üblichen logischen
Schlussregeln für \(\Vdash\). Beispielsweise gilt
\[
 p\Vdash\forall x\,\varphi(x)
 \quad\Longleftrightarrow\quad
 \text{für jeden Namen }\tau\text{ gilt }p\Vdash\varphi(\tau).
\]
Für die Hinrichtung würde eine nicht forcierte Instanz durch das
Dichtheitslemma eine Verstärkung mit forcierter Gegeninstanz
zulassen. Für die Rückrichtung würde ein forcierter existenzieller
Gegenzeuge einer der Instanzen widersprechen.
Mitgliedschaftszeugen und Gleichheitssubstitution zeigen entsprechend,
dass die atomare Hilfsrelation \(\tau\subseteq\sigma\) genau die
Formel \(\forall x(x\in\tau\to x\in\sigma)\) forciert.
Diese Ableitungen geschehen intern; sie benötigen nicht erst die
Existenz eines generischen Filters.

\FormulaThmDeltaK[Einbettung des Ausgangsmodells]
{\begin{aligned}
 j:M&\longrightarrow M[G],\qquad j(x)=\check x_G,\\
 j(x)=j(y)&\Longleftrightarrow x=y,\\
 j(x)\in j(y)&\Longleftrightarrow M\models x\in y
\end{aligned}}
{B48CheckEmbedding}{
  \DeltaRow{Abbildung}{j\text{ von außen definiert}}
}
\begin{tabproof}
 \proofstep{}{\BXLVIIIText{Interne Ranginduktion über \(x,y\)
 zeigt: Für kanonische Namen wird wahre Gleichheit und
 Mitgliedschaft von \(1\) forciert; ihre falschen Fälle
 werden von keiner Bedingung forciert.
 Die atomaren Klauseln reduzieren jeden Vergleich auf
 Elemente kleineren Ranges.}}
 {\BXLVIIIWhy{Kanonische Namen und atomare Rekursion}}
 \proofstep{}{\BXLVIIIText{Wahrheitslemma und Quotientendefinition
 übersetzen diese Aussagen in genau die beiden angegebenen
 Äquivalenzen. Insbesondere ist \(j\) injektiv.}}
 {\FormulaRefAuto{B48TruthLemma}[thm]}
\end{tabproof}
\noindent
Diese Einbettung ist keine Behauptung elementarer Äquivalenz.
Gerade die Potenzmengen und damit die Wahrheit von CH dürfen sich
ändern. Wir identifizieren alte Elemente im Folgenden mit ihren
kanonischen Bildern, ohne eine wörtliche transitive Inklusion
der äußeren Strukturen zu verlangen.

\subsection{Warum die Erweiterung wieder ZFC erfüllt}

Die folgenden Beweise behandeln jedes Axiom und jede einzelne
Instanz der Axiomenschemata. Die angegebenen Namen werden in \(M\)
gebildet. Die Dichtheits- und Zeugenargumente zeigen zugleich
intern, dass \(1\) das jeweilige Axiom forciert; das Wahrheitslemma
überträgt dies auf den generischen Quotienten.

\FormulaThmDeltaK[Extensionalität und elementare Mengenoperationen]
{M[G]\models
 \text{Extensionalität, leere Menge, Paarung und Vereinigung}}
{B48ForcingBasicAxioms}{
  \DeltaRow{Ausgangsmodell}{M\models\BXLVIIIZFC}
}
\begin{tabproof}
 \proofstep{}{\BXLVIIIText{Extensionalität: Forciert eine Bedingung
 \(\tau\ne\sigma\), so gibt es dicht darunter einen Zeugen
 für das Scheitern einer der beiden Inklusionen.
 Andernfalls würden die Inklusionsklauseln zusammen Gleichheit
 forcen. Ein solcher Zeuge gehört zu genau einer der beiden
 Mengen. Daher haben verschiedene Mengen verschiedene Elemente.}}
 {\BXLVIIIWhy{Atomare Gleichheit, Negation und Dichtheitslemma}}
 \proofstep{}{\varnothing_G\text{ hat kein Element}}
 {\FormulaRefAuto{B48NameMembers}[thm]}
 \proofstep{}{\{(\tau,1),(\sigma,1)\}_G=\{\tau_G,\sigma_G\}}
 {\FormulaRefAuto{B48NameMembers}[thm]}
 \proofstep{}{\BXLVIIIText{Für einen Namen \(\sigma\) bilde
 \[
 u=\{(\eta,q):\exists(\nu,p)\in\sigma\
               \exists(\eta,r)\in\nu\ (q\le p,r)\}.
 \]
 Dies ist eine Menge: Die möglichen Unternamen werden aus
 zwei Namensebenen gesammelt, die Bedingungen aus \(P\).}}
 {\BXLVIIIWhy{Ersetzung, Vereinigung und Aussonderung in \(M\)}}
 \proofstep{}{\BXLVIIIText{Ein Element von \(u_G\) liegt nach
 der Mitgliederbeschreibung in einem \(\nu_G\in\sigma_G\).
 Umgekehrt lassen sich die beiden dafür zuständigen Bedingungen
 aus \(G\) gemeinsam verstärken; das resultierende Paar tritt
 dann in \(u\) auf. Also \(u_G=\bigcup\sigma_G\).}}
 {\BXLVIIIWhy{Mitgliederbeschreibung und Filtergerichtetheit}}
\end{tabproof}

\FormulaThmDeltaK[Aussonderung in der Erweiterung]
{M[G]\models\text{jede feste Instanz des Aussonderungsschemas}}
{B48ForcingSeparation}{
  \DeltaRow{Formel}{\varphi(x,\bar a)}
  \DeltaRow{Namen}{\sigma,\bar\tau}
}
\begin{tabproof}
 \proofstep{}{\BXLVIIIText{In \(M\) existiert die Menge
 \[
 s=\{(\eta,q):\exists(\eta,r)\in\sigma\
       (q\le r\land q\Vdash\varphi(\eta,\bar\tau))\}.
 \]
 Die Formel ist fest, ihre Forcingrelation definierbar.}}
 {\BXLVIIIWhy{Aussonderung auf \(\operatorname{dom}(\sigma)\times P\)}}
 \proofstep{}{\BXLVIIIText{Ist \(x\in s_G\), so liefern ein Paar
 des Namens und eine Bedingung aus \(G\) sowohl
 \(x\in\sigma_G\) als auch \(\varphi(x,\bar\tau_G)\).}}
 {\BXLVIIIWhy{Mitgliederbeschreibung und Wahrheitslemma}}
 \proofstep{}{\BXLVIIIText{Sind umgekehrt \(x\in\sigma_G\) und
 \(\varphi(x,\bar\tau_G)\) wahr, so wähle einen
 Mitgliedschaftsrepräsentanten und eine die Formel forcierende
 Bedingung in \(G\). Eine gemeinsame Verstärkung gehört zum
 gerade definierten Namen \(s\).}}
 {\BXLVIIIWhy{Gleichheitssubstitution und Filtergerichtetheit}}
 \proofstep{}{s_G=\{x\in\sigma_G:\varphi(x,\bar\tau_G)\}}
 {\BXLVIIIWhy{Beide Inklusionen}}
\end{tabproof}

\FormulaThmDeltaK[Potenzmenge in der Erweiterung]
{M[G]\models\text{Potenzmengenaxiom}}
{B48ForcingPowerset}{
  \DeltaRow{Name}{\sigma}
}
\begin{tabproof}
 \proofstep{}{\BXLVIIIText{Setze in \(M\)
 \(D=\{\eta:\exists r\ ((\eta,r)\in\sigma)\}\) und
 \(R=\mathcal P(D\times P)\). Jedes Element von \(R\) ist
 ein Name. Definiere
 \(\pi=\{(\tau,p):\tau\in R,\ p\Vdash\tau\subseteq\sigma\}\).}}
 {\BXLVIIIWhy{Potenzmenge und Aussonderung in \(M\)}}
 \proofstep{}{\pi_G\subseteq\mathcal P^{M[G]}(\sigma_G)}
 {\BXLVIIIWhy{Mitgliederbeschreibung und Wahrheitslemma}}
 \proofstep{3}{\tau_G\subseteq\sigma_G}{\rA}
 \proofstep{3}{p\in G,\qquad p\Vdash\tau\subseteq\sigma}
 {\FormulaRefAuto{B48TruthLemma}[thm]}
 \proofstep{3}{\BXLVIIIText{Normalisiere den Namen unter \(p\):
 \[
 \tau'=\{(\eta,q):\exists(\eta,r)\in\sigma\
                 (q\le p,r\land q\Vdash\eta\in\tau)\}.
 \]
 Dann liegt \(\tau'\) in der festen Menge \(R\).}}
 {\BXLVIIIWhy{Aussonderung}}
 \proofstep{3}{\BXLVIIIText{Die atomaren Klauseln zeigen
 \(p\Vdash\tau'=\tau\) und \(p\Vdash\tau'\subseteq\sigma\):
 Mitglieder von \(\tau'\) sind ausdrücklich Mitglieder von
 \(\tau\) und \(\sigma\). Für ein Mitglied von \(\tau\)
 liefert die unter \(p\) forcierte Inklusion einen Vertreter
 aus \(\sigma\); eine weitere Verstärkung nimmt diesen in
 \(\tau'\) auf.}}
 {\BXLVIIIWhy{Zeugen, Substitution und Dichtheitslemma}}
 \proofstep{3}{\tau_G=\tau'_G\in\pi_G}
 {\BXLVIIIWhy{\((\tau',p)\in\pi\), \(p\in G\)}}
 \proofstep{}{\pi_G=\mathcal P^{M[G]}(\sigma_G)}
 {\BXLVIIIWhy{Beliebige Teilmenge erfasst, 3 entladen}}
\end{tabproof}

\FormulaThmDeltaK[Ersetzung in der Erweiterung]
{M[G]\models\text{jede feste Instanz des Ersetzungsschemas}}
{B48ForcingReplacement}{
  \DeltaRow{Formel}{\varphi(x,y,\bar a)}
  \DeltaRow{Name der Ausgangsmenge}{\sigma}
}
\begin{tabproof}
 \proofstep{1}{\BXLVIIIText{In \(M[G]\) ordne \(\varphi\) jedem
 \(x\in\sigma_G\) genau ein \(y\) zu. Wähle \(p\in G\),
 das diese Aussage für die entsprechenden Namen forciert.}}{\rA}
 \proofstep{1}{\BXLVIIIText{Für jedes \((\eta,r)\in\sigma\) sind
 unter gemeinsamen Verstärkungen von \(p,r\) die Bedingungen
 dicht, die \(\varphi(\eta,\nu,\bar\tau)\) für irgendeinen
 Namen \(\nu\) forcen. Denn dort ist \(\eta\in\sigma\)
 forciert und damit die Existenz des zugehörigen Wertes.}}
 {\BXLVIIIWhy{Universalklausel und Existenzklausel}}
 \proofstep{1}{\BXLVIIIText{Wähle in dieser dichten Menge eine
 maximale Antikette \(A_{\eta,r}\). Für jedes ihrer Elemente
 \(a\) wähle einen Zeugen \(\nu_{\eta,r,a}\).
 Diese Wahlen sind Mengenwahlen: Die Bedingungen liegen in
 \(P\); für jeden benötigten Zeugen wählt man zunächst den
 kleinsten möglichen Namensrang. Ersetzung begrenzt alle
 diese Ränge, danach genügt Auswahl in einer einzigen
 Namensstufe.}}
 {\BXLVIIIWhy{Auswahl und Ersetzung in \(M\);
 keine Auswahl aus einer echten Klasse}}
 \proofstep{1}{b=\{(\nu_{\eta,r,a},a):
                  (\eta,r)\in\sigma,\ a\in A_{\eta,r}\}}
 {\BXLVIIIWhy{Mengenindexierte Ersetzung und Vereinigung}}
 \proofstep{1}{\BXLVIIIText{Jedes Mitglied von \(b_G\) ist ein
 durch \(\varphi\) zugeordnetes Bild eines Mitglieds von
 \(\sigma_G\), denn die gewählte Bedingung liegt unter
 \(p,r\). Umgekehrt trifft \(G\) bei jedem
 Mitgliedschaftsvertreter die entsprechende Antikette.
 Die forcierte Eindeutigkeit identifiziert deren Zeugen mit
 dem gesuchten Wert.}}
 {\BXLVIIIWhy{Generizität, Mitgliederbeschreibung und Eindeutigkeit}}
 \proofstep{1}{b_G=\{y:\exists x\in\sigma_G\
                         \varphi(x,y,\bar\tau_G)\}}
 {\BXLVIIIWhy{Beide Inklusionen}}
\end{tabproof}

\FormulaThmDeltaK[Fundierung, Unendlichkeit und Auswahl]
{M[G]\models
 \text{Fundierung, Unendlichkeit und Auswahl}}
{B48ForcingRemainingAxioms}{
  \DeltaRow{Ausgangsmodell}{M\models\BXLVIIIZFC}
}
\begin{tabproof}
 \proofstep{}{\BXLVIIIText{Fundierung: Sei \(p\Vdash\sigma\ne\varnothing\)
 und \(q\le p\). Wähle unter allen Namen, deren Mitgliedschaft
 in \(\sigma\) irgendeine Verstärkung von \(q\) forciert,
 einen Namen \(\eta\) kleinsten Ranges; wähle dazu
 \(r\le q\) mit \(r\Vdash\eta\in\sigma\).
 Diese Minimalwahl geschieht im Ausgangsmodell.}}
 {\BXLVIIIWhy{Existenzklausel und interne Ordinalwohlordnung}}
 \proofstep{}{\BXLVIIIText{Keine Verstärkung von \(r\) kann ein
 Mitglied von \(\eta\cap\sigma\) forcen.
 Andernfalls liefert die atomare Mitgliedschaft in \(\eta\)
 nach weiterer Verstärkung einen Unternamen \(\nu\) von
 \(\eta\), der zugleich Mitglied von \(\sigma\) ist.
 Es wäre \(\rho(\nu)<\rho(\eta)\), entgegen der Minimalwahl.
 Somit forciert \(r\) ein \(\in\)-minimales Element von
 \(\sigma\). Solche \(r\) sind dicht unter \(p\).}}
 {\BXLVIIIWhy{Mitgliedschaftszeugen, Substitution und Dichtheitslemma}}
 \proofstep{}{\BXLVIIIText{Unendlichkeit: \(\check\omega_G\)
 enthält die leere Menge und mit jedem Mitglied dessen
 Nachfolger. Dies folgt aus den kanonischen Namen und der
 Mitgliederbeschreibung; sie ist also eine induktive Menge.}}
 {\FormulaRefAuto{B48CheckEmbedding}[thm]}
 \proofstep{}{\BXLVIIIText{Auswahl: Wohlordne in \(M\) den
 Namenbereich \(D=\operatorname{dom}(\sigma)\) als
 \((\eta_i)_{i<\mu}\). Ein Mengenname für den Graphen mit
 Paaren \((\check i,\eta_i)\), jeweils unter den Bedingungen,
 unter denen \(\eta_i\) als Mitglied von \(\sigma\) auftritt,
 ergibt eine Surjektion von einer Teilmenge von
 \(\check\mu_G\) auf \(\sigma_G\).}}
 {\BXLVIIIWhy{Auswahl in \(M\); Paar- und Graphennamen}}
 \proofstep{}{\BXLVIIIText{Die Menge \(\check\mu_G\) ist transitiv
 und durch \(\in\) strikt linear geordnet, nach der
 Mitgliederbeschreibung und der Einbettung.
 Die bereits bewiesene Fundierung macht sie zu einer
 Ordinalzahl. Ordne jedem Element von \(\sigma_G\) seinen
 kleinsten Urbildindex zu. So wird \(\sigma_G\) wohlgeordnet.}}
 {\BXLVIIIWhy{Aussonderung und Ersetzung in \(M[G]\)}}
 \proofstep{}{M[G]\models\mathsf{AC}}
 {\BXLVIIIWhy{Wohlordnung jeder Menge impliziert Auswahl}}
\end{tabproof}

\FormulaThmDeltaK[Erhaltung von ZFC durch Mengenforcing]
{M\models\BXLVIIIZFC,\ G\text{ \(M\)-generisch}
 \ \vdash\ M[G]\models\BXLVIIIZFC}
{B48ForcingZFC}{
  \DeltaRow{Forcingordnung}{M\models\mathbb P\text{ ist eine Menge}}
}
\begin{tabproof}
 \proofstep{1}{M\models\BXLVIIIZFC,\quad G\text{ \(M\)-generisch}}{\rA}
 \proofstep{1}{\BXLVIIIText{Extensionalität, leere Menge,
 Paarung und Vereinigung gelten durch die angegebenen
 elementaren Namen.}}
 {\FormulaRefAuto{B48ForcingBasicAxioms}[thm]}
 \proofstep{1}{\BXLVIIIText{Aussonderung und Ersetzung gelten
 für jede feste Formel; die Potenzmenge wird durch die
 Normalisierung auf einen festen Namenbereich als Menge erfasst.}}
 {\BXLVIIIWhy{Aussonderungs-, Ersetzungs- und Potenzmengenbeweis}}
 \proofstep{1}{\BXLVIIIText{Fundierung, Unendlichkeit und Auswahl
 gelten nach dem vorigen Satz. Damit ist jedes einzelne
 ZFC-Axiom im Quotienten erfüllt.}}
 {\FormulaRefAuto{B48ForcingRemainingAxioms}[thm]}
 \proofstep{1}{M[G]\models\BXLVIIIZFC}
 {\BXLVIIIWhy{Definition der Theorie ZFC}}
\end{tabproof}

\noindent
Die innere Ranginduktion wird nicht mit einer äußeren Fundiertheit
von \(M\) verwechselt. Außen benutzt das Wahrheitslemma nur
Induktion über die gewöhnliche endliche Formel.
Gerade deshalb lässt sich die Konstruktion auf die durch
Vollständigkeit gewonnenen beliebigen abzählbaren Modelle anwenden.

\subsection{Ordinale und die abzählbare Kettenbedingung}

\FormulaThmDeltaK[Keine neuen Ordinalzahlen]
{\operatorname{Ord}^{M[G]}
   =\{\check\alpha_G:M\models\alpha\in\BXLVIIIOrd\}}
{B48ForcingOrdinals}{
  \DeltaRow{Erweiterung}{M[G]}
}
\begin{tabproof}
 \proofstep{}{\BXLVIIIText{Jede alte Ordinalzahl bleibt eine
 Ordinalzahl: Transitivität und strikte \(\in\)-Linearität
 folgen aus den kanonischen Namen, Wohlgeordnetheit aus
 Fundierung in der Erweiterung.}}
 {\BXLVIIIWhy{Einbettung und Fundierung}}
 \proofstep{}{\BXLVIIIText{Interne Induktion über die Namensränge
 zeigt
 \(1\Vdash\operatorname{rk}(\tau)\le\check{\rho(\tau)}\).
 Denn jedes Mitglied von \(\tau\) wird dicht durch einen
 Unternamen \(\eta\) vertreten; dessen Rang liegt unter
 \(\rho(\tau)\). Supremum der Nachfolger der Mitgliedsränge
 liefert die Schranke.}}
 {\BXLVIIIWhy{Rangrekursion und Mitgliedschaftsklausel}}
 \proofstep{3}{p\Vdash\tau\text{ ist eine Ordinalzahl},\quad q\le p}{\rA}
 \proofstep{3}{\BXLVIIIText{Wähle intern die kleinste alte
 Ordinalzahl \(\alpha\le\rho(\tau)\), für die
 \(q\not\Vdash\check\alpha\in\tau\).
 Eine solche gibt es wegen der Rangschranke und
 \(\operatorname{rk}(\tau)=\tau\) für Ordinale.}}
 {\BXLVIIIWhy{Interne Wohlordnung eines Mengenintervalls}}
 \proofstep{3}{\BXLVIIIText{Nach dem Dichtheitslemma gibt es
 \(r\le q\) mit \(r\Vdash\check\alpha\notin\tau\).
 Für jedes \(\beta<\alpha\) forciert schon \(q\) die
 Mitgliedschaft von \(\check\beta\) in \(\tau\).
 Die Mitgliederbeschreibung kanonischer Namen macht daraus
 \(r\Vdash\check\alpha\subseteq\tau\).}}
 {\BXLVIIIWhy{Minimalität; jedes Mitglied von \(\check\alpha\)
 wird dicht durch ein altes \(\check\beta\) vertreten}}
 \proofstep{3}{r\Vdash\tau=\check\alpha}
 {\BXLVIIIWhy{Ordinalvergleich; echte Obermenge würde
 \(\check\alpha\in\tau\) erzwingen}}
 \proofstep{}{\BXLVIIIText{Die Bedingungen, welche den Wert
 eines Ordinalnamens auf eine alte Ordinalzahl festlegen,
 sind daher dicht unter jeder Bedingung, die sein
 Ordinalsein forciert. Generizität und Wahrheitslemma
 liefern die Behauptung.}}
 {\BXLVIIIWhy{3 entladen}}
\end{tabproof}
\noindent
Insbesondere ist \(\omega^{M[G]}=\check\omega_G\):
Die alte \(\omega\) bleibt eine nichtnullige Limesordinalzahl,
und alle kleineren alten Ordinale bleiben \(0\) oder Nachfolger.
Das Argument gilt intern auch dann, wenn \(M\) von außen
nichtstandardmäßige natürliche Zahlen besitzt.

\FormulaDefDeltaK[Abzählbare Kettenbedingung]
{\mathbb P\text{ erfüllt ccc}\ \Longleftrightarrow\
 \forall A\subseteq P\
 (A\text{ Antikette}\ \longrightarrow\ |A|\le\aleph_0)}
{B48CCCDef}{
  \DeltaRow{Abkürzung}{\text{ccc: countable chain condition}}
}

\FormulaThmDeltaK[ccc erhält die Kardinalzahlen]
{M\models\text{\(\mathbb P\) erfüllt ccc}
 \ \vdash\ \text{\(M[G]\) und \(M\) haben dieselben Kardinalzahlen}}
{B48CCCPreservesCardinals}{
  \DeltaRow{Identifikation}{\alpha\longmapsto\check\alpha_G}
}
\begin{tabproof}
 \proofstep{1}{M\models\text{\(\mathbb P\) erfüllt ccc}}{\rA}
 \proofstep{2}{\BXLVIIIText{Eine alte unendliche Kardinalzahl
 \(\kappa\) sei in der Erweiterung keine Kardinalzahl mehr.
 Dann gibt es eine alte Ordinalzahl \(\mu<\kappa\) und eine
 Surjektion \(f:\mu\to\kappa\). Wähle \(p\in G\), das
 diese Aussage für einen Namen von \(f\) forciert.}}{\rA}
 \proofstep{1,2}{\BXLVIIIText{Für jedes \(\xi<\mu\) wähle in
 \(M\) eine maximale Antikette unter \(p\), deren Bedingungen
 den Ordinalwert \(f(\xi)\) entscheiden. Solche Entscheider
 sind dicht, weil keine neuen Ordinale entstehen.
 Jede Antikette ist nach ccc abzählbar; daher liegt
 \(f(\xi)\) in einer alten abzählbaren Menge \(B_\xi\subseteq\kappa\).}}
 {\BXLVIIIWhy{Entscheidung von Ordinalnamen; Auswahl in \(M\)}}
 \proofstep{1,2}{\BXLVIIIText{Ist \(\kappa>\aleph_0\), so gilt
 in \(M\)
 \(\bigl|\bigcup_{\xi<\mu}B_\xi\bigr|
 \le\max(|\mu|,\aleph_0)<\kappa\).
 Wähle dort ein \(\beta\) außerhalb dieser Vereinigung.
 Die maximalen Antiketten erzwingen, dass \(f\) dieses
 \(\beta\) nicht trifft, im Widerspruch zur Surjektivität.}}
 {\FormulaRefAuto{B48CardinalProduct}[thm]}
 \proofstep{1,2}{\BXLVIIIText{Für \(\kappa=\aleph_0\) müsste
 eine endliche Menge auf \(\omega\) surjizieren, was in ZFC
 unmöglich ist. Die natürlichen Zahlen und \(\omega\)
 sind unter der kanonischen Einbettung unverändert.}}
 {\BXLVIIIWhy{Endliche Arithmetik und Erhaltung von \(\omega\)}}
 \proofstep{1}{\BXLVIIIText{Keine alte Kardinalzahl wird also
 kollabiert. Alte Nichtkardinale bleiben Nichtkardinale,
 denn eine alte Bijektion mit einer kleineren Ordinalzahl
 bleibt bestehen. Neue Ordinalzahlen gibt es nicht.}}
 {\BXLVIIIWhy{2 entladen; Einbettung und Ordinalerhaltung}}
\end{tabproof}

\subsection{Cohen-Bedingungen und neue reelle Zahlen}

\FormulaThmDeltaK[Delta-System-Lemma für endliche Mengen]
{\begin{gathered}
 \mathcal A\text{ überabzählbare Familie verschiedener endlicher Mengen}\\
 \vdash\ \exists\mathcal B\subseteq\mathcal A\
 \exists r\quad |\mathcal B|=\aleph_1,\quad
 \forall a\ne b\in\mathcal B\ (a\cap b=r)
 \end{gathered}}
{B48DeltaSystem}{
  \DeltaRow{Menge}{\mathcal A}
}
\begin{tabproof}
 \proofstep{1}{\mathcal A\text{ wie vorausgesetzt}}{\rA}
 \proofstep{1}{\BXLVIIIText{Beschränke auf eine Familie der
 Mächtigkeit \(\aleph_1\). Für eine feste natürliche Zahl
 \(n\) gibt es überabzählbar viele Mitglieder mit genau
 \(n\) Elementen; sonst wäre die ganze Familie abzählbar.
 Induktion über \(n\).}}
 {\BXLVIIIWhy{Abzählbare Vereinigung abzählbarer Mengen}}
 \proofstep{1}{\BXLVIIIText{Kommt ein Element \(x\) in
 überabzählbar vielen Mitgliedern vor, entferne \(x\)
 aus diesen Mitgliedern. Sie bleiben verschieden und
 haben Größe \(n-1\). Die Induktionsvoraussetzung liefert
 ein Delta-System; füge \(x\) zu seiner Wurzel hinzu.}}
 {\BXLVIIIWhy{Induktionsschritt, erster Fall}}
 \proofstep{1}{\BXLVIIIText{Andernfalls kommt jedes Element
 nur in abzählbar vielen Mitgliedern vor.
 Wähle rekursiv \(\aleph_1\) viele paarweise disjunkte
 Mitglieder: Vor jeder Stufe \(\xi<\omega_1\) wurden nur
 abzählbar viele endliche Mengen gewählt; ihre Vereinigung
 trifft nur abzählbar viele verbotene Mitglieder.
 Es bleibt ein weiteres Mitglied übrig.}}
 {\BXLVIIIWhy{Induktionsschritt, zweiter Fall; Wurzel \(\varnothing\)}}
 \proofstep{1}{\BXLVIIIText{Der Anfang \(n=0\) ist leer, da
 es nicht überabzählbar viele verschiedene leere Mengen
 gibt. Beide Fälle schließen die Induktion ab.}}
 {\BXLVIIIWhy{Induktion auf den natürlichen Zahlen}}
\end{tabproof}

\FormulaDefDeltaK[Cohen-Forcing für \(\lambda\) reelle Zahlen]
{\begin{aligned}
 \operatorname{Add}(\omega,\lambda)
 &=\{p:p\text{ endliche partielle Funktion }
                         \lambda\times\omega\longrightarrow2\},\\
 p\le q&\ \Longleftrightarrow\ p\supseteq q,\qquad 1=\varnothing
\end{aligned}}
{B48CohenOrder}{
  \DeltaRow{Unendliche Kardinalzahl}{\lambda}
}

\FormulaThmDeltaK[Größe und ccc des Cohen-Forcings]
{\lambda\ge\aleph_0\ \vdash\
 |\operatorname{Add}(\omega,\lambda)|=\lambda
 \quad\land\quad \operatorname{Add}(\omega,\lambda)\text{ erfüllt ccc}}
{B48CohenCCC}{
  \DeltaRow{Kardinalzahl}{\lambda}
}
\begin{tabproof}
 \proofstep{1}{\lambda\ge\aleph_0}{\rA}
 \proofstep{1}{\BXLVIIIText{Eine Bedingung wird durch endlich
 viele Tripel aus \(\lambda\times\omega\times2\) codiert.
 Daher gibt es höchstens \(\lambda\) Bedingungen.
 Die Einzelbedingungen an den Stellen \((\alpha,0)\)
 liefern die untere Schranke \(\lambda\).}}
 {\FormulaRefAuto{B48CardinalProduct}[thm]}
 \proofstep{3}{\mathcal A\text{ überabzählbare Antikette}}{\rA}
 \proofstep{1,3}{\BXLVIIIText{Es treten überabzählbar viele
 verschiedene endliche Definitionsbereiche auf, da jeder
 feste endliche Bereich nur endlich viele Funktionen nach
 \(2\) zulässt. Wähle je einen Vertreter und ein
 überabzählbares Delta-System seiner Bereiche mit Wurzel \(r\).}}
 {\FormulaRefAuto{B48DeltaSystem}[thm]}
 \proofstep{1,3}{\BXLVIIIText{Auf der endlichen Wurzel gibt es
 nur endlich viele mögliche Belegungen. Zwei der gewählten
 Bedingungen, sogar überabzählbar viele, stimmen dort
 überein. Außerhalb der Wurzel sind ihre Bereiche disjunkt;
 ihre Vereinigung ist daher eine gemeinsame Verstärkung.}}
 {\BXLVIIIWhy{Endliches Schubfachprinzip und Funktionsvereinigung}}
 \proofstep{1}{\operatorname{Add}(\omega,\lambda)\text{ erfüllt ccc}}
 {\BXLVIIIWhy{Widerspruch zur Antikette, 3 entladen}}
\end{tabproof}

\FormulaThmDeltaK[Mindestens \(\lambda\) verschiedene neue Binärfolgen]
{\begin{gathered}
 M\models\lambda\ge\aleph_0,\quad
 \mathbb P=\operatorname{Add}(\omega,\lambda)^M\\
 \vdash\ M[G]\models \lambda\le2^{\aleph_0}
 \end{gathered}}
{B48CohenReals}{
  \DeltaRow{Filter}{G\text{ \(M\)-generisch}}
}
\begin{tabproof}
 \proofstep{}{\BXLVIIIText{Für jede alte Stelle \((\alpha,n)\)
 ist die Menge der Bedingungen, die dort einen Wert
 festlegen, dicht: Fehlt die Stelle, füge einen der Werte
 \(0,1\) hinzu. Vereinbarkeit in einem Filter verhindert
 widersprüchliche Werte.}}
 {\BXLVIIIWhy{Endliche partielle Funktionen}}
 \proofstep{}{\BXLVIIIText{Für verschiedene \(\alpha,\beta\)
 ist auch die Menge der Bedingungen dicht, die an einer
 Stelle \(n\) verschiedene Werte für \((\alpha,n)\) und
 \((\beta,n)\) festlegen. Man wählt ein von den endlich
 vielen belegten Stellen noch nicht betroffenes \(n\).}}
 {\BXLVIIIWhy{Endlichkeit der Bedingungen}}
 \proofstep{}{\BXLVIIIText{Die Namen
 \(\dot c_\alpha=\{(\check n,p):p(\alpha,n)=1\}\)
 definieren daher paarweise verschiedene Teilmengen von
 \(\omega\). Die ganze Familie gehört zur Erweiterung:
 Ein Mengenname für ihren Graphen wird aus
 \((\langle\check\alpha,\dot c_\alpha\rangle,1)\),
 \(\alpha<\lambda\), gebildet.}}
 {\BXLVIIIWhy{Ersetzung in \(M\); Paar- und Graphennamen}}
 \proofstep{}{\BXLVIIIText{Fixiere \(\alpha<\lambda\) und eine
 alte Binärfolge \(x\in(2^\omega)^M\). Die in \(M\) definierte
 Menge \(D_{\alpha,x}\) aller Bedingungen, die an einer Stelle
 \((\alpha,n)\) den Wert \(1-x(n)\) festlegen, ist dicht:
 Zu jeder Bedingung wähle ein noch unbelegtes \(n\) und füge
 genau diesen Wert hinzu.}}
 {\BXLVIIIWhy{Endlichkeit; Aussonderung in \(M\)}}
 \proofstep{}{\BXLVIIIText{Da \(G\) die Menge
 \(D_{\alpha,x}\) trifft, unterscheidet sich die generische
 Binärfolge an dieser Stelle von \(x\). Dies gilt für jede
 alte Binärfolge. Über charakteristische Funktionen ist
 somit auch \(\dot c_\alpha^G\) keine alte Teilmenge von
 \(\omega\). Die Folgen sind also nicht nur untereinander
 verschieden, sondern tatsächlich neu.}}
 {\BXLVIIIWhy{Generizität; kanonische Einbettung von \(M\)}}
 \proofstep{}{M[G]\models\lambda\hookrightarrow\mathcal P(\omega)}
 {\BXLVIIIWhy{Generizität für die angegebenen dichten Mengen}}
 \proofstep{}{M[G]\models\lambda\le2^{\aleph_0}}
 {\BXLVIIIWhy{ccc erhält \(\lambda,\omega\); charakteristische Funktionen}}
\end{tabproof}

\subsection{Zählen von Namen: die notwendige obere Schranke}

\FormulaThmDeltaK[Schöne Namen für Teilmengen einer Ordinalzahl]
{\begin{gathered}
 x\in M[G],\quad M[G]\models x\subseteq\mu\\
 \vdash\ x=\nu_G,\qquad
 \nu=\{(\check\xi,p):\xi<\mu,\ p\in A_\xi\},\\
 A_\xi\text{ ist eine Antikette in }\mathbb P\text{ aus }M
 \end{gathered}}
{B48NiceNames}{
  \DeltaRow{Alte Ordinalzahl}{\mu}
}
\begin{tabproof}
 \proofstep{1}{x=\tau_G\subseteq\mu}{\rA}
 \proofstep{1}{\BXLVIIIText{Wähle in \(M\) für jedes \(\xi<\mu\)
 eine maximale Antikette \(A_\xi\) in der Menge
 \(\{p:p\Vdash\check\xi\in\tau\}\).
 Diese Menge kann leer sein; dann ist die Antikette leer.
 Ersetzung und Auswahl sammeln die Antiketten zu einer Folge.}}
 {\BXLVIIIWhy{Definierbarkeit des atomaren Forcings; Auswahl}}
 \proofstep{1}{\nu=\bigcup_{\xi<\mu}
                \{(\check\xi,p):p\in A_\xi\}}
 {\BXLVIIIWhy{Mengenbildung in \(M\)}}
 \proofstep{1}{\BXLVIIIText{Trifft \(G\) die Antikette \(A_\xi\),
 so gilt \(\xi\in\tau_G\). Gilt umgekehrt
 \(\xi\in\tau_G\), so gibt es \(q\in G\), das diese
 Mitgliedschaft forciert. Unter \(q\) sind die
 Verstärkungen der Elemente von \(A_\xi\) dicht:
 Jede Verstärkung von \(q\) gehört zur positiven
 Forcingmenge und ist nach Maximalität mit einem
 Antikettenelement verträglich.}}
 {\BXLVIIIWhy{Wahrheitslemma und Maximalität}}
 \proofstep{1}{\BXLVIIIText{Generizität unter \(q\) liefert
 deshalb ein Element von \(A_\xi\) in \(G\).
 Somit haben \(\nu_G\) und \(\tau_G\) genau dieselben
 Mitglieder aus \(\mu\), und nach Voraussetzung keine
 weiteren.}}
 {\BXLVIIIWhy{Mitgliederbeschreibung; \(x\subseteq\mu\)}}
 \proofstep{1}{\nu_G=x}{\BXLVIIIWhy{Extensionalität}}
\end{tabproof}

\FormulaThmDeltaK[Zählschranke bei ccc]
{\begin{gathered}
 M\models\bigl(\mathbb P\text{ ccc},\
 |\mathbb P|=\lambda\ge\aleph_0,\ \mu\ge\aleph_0\bigr)\\
 K=\bigl((\lambda^{\aleph_0})^\mu\bigr)^M
 \quad\Longrightarrow\quad M[G]\models 2^\mu\le K
 \end{gathered}}
{B48NiceNameBound}{
  \DeltaRow{Alte Kardinale}{\lambda,\mu,K}
}
\begin{tabproof}
 \proofstep{1}{M\models\mathbb P\text{ ccc},\quad |P|=\lambda\ge\aleph_0}
 {\rA}
 \proofstep{1}{\BXLVIIIText{Jede Antikette ist in \(M\)
 abzählbar. Die nichtleeren abzählbaren Teilmengen von
 \(P\) werden von Folgen in \(P^\omega\) aufgezählt;
 Hinzunahme der leeren Menge ändert die unendliche
 obere Schranke \(\lambda^{\aleph_0}\) nicht.}}
 {\BXLVIIIWhy{ccc und Auswahl in \(M\)}}
 \proofstep{1}{\BXLVIIIText{Ein schöner Name wird durch eine
 \(\mu\)-Folge solcher Antiketten bestimmt.
 Die Menge \(\mathcal N\) aller schönen Namen hat daher
 in \(M\) Mächtigkeit höchstens
 \(K=(\lambda^{\aleph_0})^\mu\).}}
 {\BXLVIIIWhy{Funktionszählung}}
 \proofstep{1}{\BXLVIIIText{Wähle eine Aufzählung
 \((\nu_i)_{i<K}\) aller dieser Namen, gegebenenfalls mit
 Wiederholungen. Der Graphenname mit Paaren
 \((\check i,\nu_i)\) liefert in der Erweiterung eine
 Abbildung auf sämtliche Teilmengen von \(\mu\):
 Jede solche Teilmenge besitzt nach dem vorigen Satz
 einen schönen Namen.}}
 {\FormulaRefAuto{B48NiceNames}[thm]}
 \proofstep{1}{M[G]\models|\mathcal P(\mu)|\le K}
 {\BXLVIIIWhy{Surjektion, Auswahl in \(M[G]\);
 ccc erhält die alte Kardinalzahl \(K\)}}
\end{tabproof}

\noindent
Alle Zählungen finden zunächst \emph{in \(M\)} statt.
Die äußere Abzählbarkeit von \(M\) sagt nichts darüber aus,
welche Mächtigkeiten \(M\) oder \(M[G]\) intern besitzen.
Die Mengen von Namen und ihre Aufzählungsgraphen sind deshalb
ausdrücklich innerhalb des Ausgangsmodells konstruiert worden.
